\begin{trapbox}
  \textbf{\textcolor{CTrap}{易错 1｜把“点左减”和“函数左加”记反}}

  \textbf{错误：}$y=x^2$ 向左移 $2$ 格，写成 $y=(x-2)^2$。

  \textbf{正确：}左移后顶点是 $(-2,0)$，所以 $y=(x+2)^2$。\textbf{点的横坐标减少 $2$，新解析式中的 $x$ 却要加 $2$。}
  \begin{center}
  \begin{tikzpicture}[x=0.87cm,y=0.63cm,>=stealth,font=\small]
    \draw[->,gray] (-3.4,0)--(2.2,0) node[right] {$x$};
    \draw[->,gray] (0,-0.28)--(0,2.4) node[above] {$y$};
    \draw[blue!75!black,thick,domain=-1.5:1.5,samples=45] plot (\x,{\x*\x});
    \draw[teal!75!black,thick,domain=-3.45:-0.55,samples=45] plot (\x,{(\x+2)*(\x+2)});
    \fill[blue!75!black] (0,0) circle (2.2pt);
    \fill[teal!75!black] (-2,0) circle (2.2pt);
    \node[below] at (-2,0) {$-2$};
    \draw[->,teal!75!black,thick] (-0.25,-0.58)--(-1.75,-0.58)
      node[midway,below] {左移};
    \node[teal!75!black,above] at (-2.2,2.1) {$y=(x+2)^2$};
  \end{tikzpicture}

  \end{center}

  \medskip
  \textbf{\textcolor{CTrap}{易错 2｜只在式子末尾加减，没换掉整个 $x$}}

  \textbf{错误：}$f(x)=2x$ 左移 $1$，写成 $g(x)=2x+1$。

  \textbf{正确：}把\textbf{原式中的 $x$} 换成 $x+1$：
  \[
    g(x)=f(x+1)=2(x+1)=\boxed{2x+2}.
  \]
  \textbf{怎么检查：}原图过 $(0,0)$，左移后必须过 $(-1,0)$。$2x+2$ 在 $x=-1$ 时为 $0$，$2x+1$ 却不是。
  \begin{center}
  \begin{tikzpicture}[x=1.05cm,y=0.48cm,>=stealth,font=\small]
    \draw[->,gray] (-2.1,0)--(2.2,0) node[right] {$x$};
    \draw[->,gray] (0,-2.9)--(0,3.2) node[above] {$y$};
    \draw[thick,blue!75!black,domain=-1.4:1.45,samples=2] plot (\x,{2*\x});
    \draw[very thick,teal!75!black,domain=-2:0.5,samples=2] plot (\x,{2*\x+2});
    \draw[densely dashed,orange!85!black,thick,domain=-1.8:0.9,samples=2] plot (\x,{2*\x+1});
    \fill[blue!75!black] (0,0) circle (2.2pt);
    \fill[teal!75!black] (-1,0) circle (2.2pt);
    \node[below] at (-1,0) {$-1$};
    \node[teal!75!black,above] at (-1.25,1.6) {$2x+2$（正确）};
    \node[orange!85!black,right] at (2.2,4.0) {$2x+1$（错）};
    \node[blue!75!black,right] at (2.2,2.4) {$2x$};
  \end{tikzpicture}

  \end{center}
  \textbf{再提醒：}若 $f(x)=\sin(2x)$ 向右移 $1$，应写 $\sin\!\bigl(2(x-1)\bigr)$，不是 $\sin(2x-1)$；\textbf{先替换 $x$，再化简}。

  \medskip
  \textbf{\textcolor{CTrap}{易错 3｜只改了式子，忘了定义域也要移动}}

  \textbf{题目：}$y=\sqrt{x-1}$ 向右移 $1$ 格。

  新式为 $y=\sqrt{x-2}$，原来 $x\geqslant1$，现在\textbf{$x\geqslant2$}。
  不能沿用原来的 $x\geqslant1$，也不用担心图像“超出了原定义域”：它就是\textbf{整体搬到了新位置}。
  \begin{center}
  \begin{tikzpicture}[x=1.12cm,y=0.75cm,>=stealth,font=\small]
    \draw[->,gray] (0.3,0)--(5.55,0) node[right] {$x$};
    \draw[very thick,blue!75!black,->] (1,0.85)--(5.18,0.85);
    \draw[very thick,teal!75!black,->] (2,0.3)--(5.18,0.3);
    \fill[blue!75!black] (1,0.85) circle (2pt);
    \fill[teal!75!black] (2,0.3) circle (2pt);
    \draw[densely dashed,gray] (1,0.05)--(1,1.05);
    \draw[densely dashed,gray] (2,0.05)--(2,1.05);
    \node[below] at (1,0) {$1$};
    \node[below] at (2,0) {$2$};
    \node[above,blue!75!black] at (3.1,0.85) {原定义域 $x\geqslant1$};
    \node[below,teal!75!black] at (3.3,0.3) {新定义域 $x\geqslant2$};
    \draw[->,teal!75!black,thick] (1.15,1.45)--(1.85,1.45)
      node[midway,above] {右移 $1$};
  \end{tikzpicture}

  \end{center}

  \textbf{\textcolor{CTrap}{一招自检：}}找一个好认的点，检查它是否按题目方向移动了指定距离，\textbf{纵坐标是否不变}。
\end{trapbox}
